%%% FIELD proof-prop-known-reductions
First suppose that \(M_1=M_2=:M\), and fix a value \(\kappa\in\Theta(P)\) for which the displayed common-menu premise holds.  For \(d,d'\in M\), define
\[
 g_{1d}:=-p_{1d}-\mathbf 1\{h_0\ne d\}\kappa,
 \qquad
 g_{2d}(d'):=-p_{2d}-\mathbf 1\{d'\ne d\}\kappa,
 \qquad
 \Delta_d(d'):=g_{2d}(d')-g_{1d}.
\]
These indices are obtained from the utility index by subtracting the common constant \(\kappa\) at each date.  The premise is precisely
\[
 \min_{d\in A}\Delta_d(d')\ge
 \max_{c\in M\setminus A}\Delta_c(d')
 \quad\text{for every }d'\in A.
\]
Because \(A\) and \(M\setminus A\) are nonempty, this is equivalent to
\[
 \Delta_d(d')\ge \Delta_c(d')
 \quad
 \text{for every }d,d'\in A\text{ and }c\in M\setminus A.
\tag{1}
\]

Let \(f\in F_\kappa\), and choose a point \(z\) in its strict response patch, whose existence follows from the definition of \(F_\kappa\).  Suppose that \(f_1=d\in A\).  The strict first-period winner inequalities in the comparison graph give, for every \(c\in M\setminus A\),
\[
 g_{1d}+z_d-(g_{1c}+z_c)>0.
\tag{2}
\]
For any counterfactual lag \(d'\in A\), equations (1) and (2) imply
\[
 \begin{aligned}
 g_{2d}(d')+z_d-[g_{2c}(d')+z_c]
 &=g_{1d}+z_d-(g_{1c}+z_c)
   +\Delta_d(d')-\Delta_c(d')\\
 &>0
 \end{aligned}
\tag{3}
\]
for every \(c\notin A\).  Hence no alternative outside \(A\) can be the period-two winner under lag \(d'\), and therefore
\[
 \mathbf 1\{f_1\in A\}\le \mathbf 1\{f_2(d')\in A\}
 \quad\text{for every }d'\in A.
\tag{4}
\]
This is the samplewise containment furnished by the displayed dominance condition.

For \(q\in\Delta(F_\kappa)\), put
\[
 x(q):=\sum_{a\in A}r_a(q),
 \qquad
 u_a(q):=\sum_{b\in A}s_{ab}(q).
\]
Averaging (4) with respect to \(q\) yields \(u_a(q)\ge x(q)\) for every \(a\in A\).  Since every response table selects exactly one second-period alternative after each history,
\(\sum_{b\in M}s_{ab}(q)=1\).  The definition of the quadratic map consequently gives
\[
 \sum_{a\in A}\sum_{b\in A}v_{ab}(q)
 =\sum_{a\in A}r_a(q)u_a(q)
 \ge x(q)^2,
 \qquad
 \sum_{a\in A}\sum_{b\in M}v_{ab}(q)=x(q).
\tag{5}
\]

By the definition of \(\Theta(P)\), \(P\in\mathcal C_\kappa\); the sharp conditional-iid body theorem establishes that this membership is exactly the maintained-model restriction.  By the definition of the convex hull there are \(q^1,\ldots,q^L\in\Delta(F_\kappa)\) and weights \(\omega_\ell\ge0\), \(\sum_{\ell=1}^L\omega_\ell=1\), such that
\(P=\sum_{\ell=1}^L\omega_\ell v(q^\ell)\).  Write \(x_\ell=x(q^\ell)\).  Applying (5) componentwise and using the elementary identity
\[
 \sum_{\ell=1}^L\omega_\ell x_\ell^2-
 \left(\sum_{\ell=1}^L\omega_\ell x_\ell\right)^2
 =\frac12\sum_{\ell,j=1}^L
   \omega_\ell\omega_j(x_\ell-x_j)^2\ge0
\]
proves
\[
 \Pr(Y_1\in A,Y_2\in A\mid Y_0=h_0)
 \ge
 \Pr(Y_1\in A\mid Y_0=h_0)^2.
\tag{6}
\]
Thus (6) derives, within the present sharp-body representation, the restriction attributed in the statement to Pakes et al.  If
\(x:=\Pr(Y_1\in A\mid Y_0=h_0)>0\), division of (6) by \(x\) is valid and gives the displayed conditional-probability ratio; multiplication by \(x\) proves the converse equivalence.  If \(x=0\), that ratio is not determined by the observed law, so division is invalid.  If one formally takes \(A=M\), both probabilities in (6) equal one and the inequality is the tautology \(1\ge1\).

The same argument cannot be applied directly to a changing full menu.  Indeed, its algebra requires both \(g_{1d}\) and \(g_{2d}(d')\), and hence both period-specific prices, for every inside alternative \(d\) and every outside alternative \(c\) in one common partition.  When \(M_1\ne M_2\), an entrant has no first-period comparison and an exiting plan has no second-period comparison.  Thus the common-menu difference \(\Delta_d(d')\) and the equality in (3) are not simultaneously defined over the full changing choice universe.  This proves only the claimed inapplicability of that common-menu condition, not the failure of every possible inequality under menu turnover.

We next establish the response-type reduction.  For a fixed response table \(f\), write the finitely many simple-cycle weights as
\[
 L_C(\kappa)=A_C+B_C\kappa.
\]
The strict-cycle feasibility lemma and the definition of \(I_f\) give
\[
 f\in F_\kappa
 \quad\Longleftrightarrow\quad
 L_C(\kappa)>0\text{ for every simple cycle }C
 \quad\Longleftrightarrow\quad
 \kappa\in I_f.
\tag{7}
\]
By the definition of \(F_\infty\), \(f\in F_\infty\) exactly when every cycle has either \(B_C>0\), or \(B_C=0\) and \(A_C>0\), and no cycle has \(B_C<0\).  In that event, finiteness of the cycle list permits choosing
\[
 k_f>\max\left(\{0\}\cup
 \{-A_C/B_C:B_C>0\}\right),
\]
with the maximum over an empty second set understood as zero.  Then \((k_f,+\infty)\subseteq I_f\).  Conversely, if such a terminal interval is contained in \(I_f\), a coefficient \(B_C<0\) would make \(L_C(\kappa)<0\) for all sufficiently large \(\kappa\), and a coefficient \(B_C=0\) would require \(A_C>0\).  Hence
\[
 f\in F_\infty
 \quad\Longleftrightarrow\quad
 (k_f,+\infty)\subseteq I_f\text{ for some finite }k_f.
\tag{8}
\]

Moreover, if \(I_f\ne\varnothing\) is unbounded above, no inequality defining it can have \(B_C<0\), and every inequality with \(B_C=0\) must have \(A_C>0\).  All remaining inequalities have positive slope and hence hold on a common terminal interval.  Thus, for nonempty \(I_f\), unboundedness is equivalent to the terminal-interval property in (8).  Combining this observation with (7)--(8),
\[
 \bigl[(\exists\kappa\ge0:f\in F_\kappa)\Longrightarrow f\in F_\infty\bigr]
 \quad\Longleftrightarrow\quad
 [I_f=\varnothing\text{ or }I_f\text{ is unbounded above}].
\tag{9}
\]
Taking (9) over all response tables proves the asserted design-by-design criterion.

Under that criterion, \(F_\kappa\subseteq F_\infty\) for every finite \(\kappa\).  A probability vector on \(F_\kappa\) extends to one on \(F_\infty\) by assigning zero mass to the additional tables, without changing \(r\), \(s\), or \(v\).  The compatible-body definitions therefore imply
\[
 \mathcal C_\kappa\subseteq\mathcal C_\infty
 \quad\text{for every }\kappa\ge0.
\tag{10}
\]
Correct specification supplies \(\kappa_0\in\Theta(P)\), so \(P\in\mathcal C_{\kappa_0}\subseteq\mathcal C_\infty\).  The terminal-frontier theorem gives \(\mathcal C_\kappa=\mathcal C_\infty\) for all \(\kappa>K\).  Hence every \(\kappa>K\) belongs to \(\Theta(P)\), and the definition of the upper frontier yields \(U(P)=+\infty\).  Notice that this argument uses terminal equality and does not assume monotonicity of \(\mathcal C_\kappa\).

Finally consider the stated pure-entry design.  Let
\[
 x:=z_a-z_b,
 \qquad
 \delta_1:=p_{1a}-p_{1b},
 \qquad
 \delta_2:=p_{2a}-p_{2b}.
\]
For the response table with \(f_1=b\), \(f_2(a)=a\), and \(f_2(b)=b\), the comparison-graph definition gives, respectively,
\[
 x<\delta_1-\kappa,
 \qquad
 \delta_2-\kappa<x,
 \qquad
 x<\delta_2+\kappa.
\tag{11}
\]
The open interval in (11) is nonempty exactly when
\[
 \delta_2-\kappa<\delta_1-\kappa
 \quad\text{and}\quad
 \delta_2-\kappa<\delta_2+\kappa,
\]
which are, respectively, \(\tau=\delta_1-\delta_2>0\) and \(2\kappa>0\).  Since the maintained parameter space has \(\kappa\ge0\), the latter condition is equivalent to \(\kappa>0\).  These are the two opposing-cycle conditions identified by the strict-cycle lemma.

They are also sufficient for the full three-plan system.  Choose any \(x\) in the interval (11), set \(z_b=0\) and \(z_a=x\), and then choose \(z_c\) so small that
\[
 z_c-z_a<p_{2c}-p_{2a}+\kappa,
 \qquad
 z_c-z_b<p_{2c}-p_{2b}+\kappa.
\tag{12}
\]
There is no first-period comparison involving \(c\), because \(c\notin M_1\), and (12) satisfies both comparisons against the never-chosen entrant.  Hence all strict graph constraints hold.  We have proved
\[
 f\in F_\kappa
 \quad\Longleftrightarrow\quad
 \tau>0\text{ and }\kappa>0.
\]
In particular, \(\tau=0\) is infeasible, while for \(\tau>0\) the threshold zero is approached but not attained.

For completeness, the singleton boundaries follow directly from the same definitions.  If \(M_1=\{a\}\), period one is forced, each possible period-two winner can be made strict by raising its \(z\)-coordinate, and the quadratic map is \(v_{ab}(q)=q_b\); hence the compatible body is the simplex on \(M_2\).  If \(M_2=\{b\}\), period two is forced, each possible period-one winner can similarly be made strict, and \(v_{ab}(q)=q_a\); hence the body is the simplex on \(M_1\).  If both menus are singletons, the comparison graph is empty and these simplices reduce to the unique point mass.  This includes all stated degenerate boundary cases.
